\section{官方幻灯片证据链：Data、Grader 与 Efficiency}

本节沿官方 deck 重建工业 Agent 训练的主线。讲者没有把成功归因给单个 RL loss，而是反复强调三个互相制约的系统：数据决定模型能探索什么，grader 决定什么行为会被强化，基础设施决定单位时间能获得多少有效 rollout。读图时应把算法曲线放回这三项约束中。

\subsection{训练目标：从 coding simulator 看完整 Agent 闭环}

slides 2--4 先交代讲者负责 Microsoft CoreAI post-training 的背景，再展示 coding-agent simulator：problem prompt、code repository、language model、tools、Linux system 与 grader 共同构成环境。Agentic training 同时包含 goal orientation、planning/reasoning、tool usage 和 user interaction；任何一项脱离环境都不能代表端到端能力。

\lecturefigure{slides-images/slide-002.png}{官方 slide 2：讲者与 CoreAI post-training 背景}
\lecturefigure{slides-images/slide-003.png}{官方 slide 3：Coding Agent environment simulator}
\lecturefigure{slides-images/slide-004.png}{官方 slide 4：Goal、planning/reasoning、tool usage 与 user interaction}

\begin{knowledgebox}{读图：训练对象是闭环，不是模型输出}
prompt 与 repository 构成初始状态，模型生成动作，工具和 Linux 环境返回新观察，grader 检查最终结果。若训练数据只保留最终 patch，模型学不到何时搜索、运行测试或回退；若只评 unit test，又可能忽略用户约束。完整轨迹和多层 grader 才能把闭环转换成学习信号。
\end{knowledgebox}

\subsection{Data：可验证、Rubric、Curriculum 与 Synthesis}

slides 6--13 从数据类型走到合成流水线。数学/代码等 verifiable task 可用程序奖励，开放任务需要 detailed rubric。slide 8 追问是否一个样本也能训练，slide 9 用历史训练准确率方差对数据排序，说明 curriculum 的信息量比绝对数量更重要。高质量 data mix 往往胜过参数调优“炼金术”；Agentic synthesis 则先构造可 build 的 repository/image，再生成有 verifier 的任务。

\lecturefigure{slides-images/slide-006.png}{官方 slide 6：Verifiable 与 non-verifiable RL data}
\lecturefigure{slides-images/slide-007.png}{官方 slide 7：可解释、可 steer 的 detailed rubric}
\lecturefigure{slides-images/slide-008.png}{官方 slide 8：一个数据点能否产生训练信号}
\lecturefigure{slides-images/slide-009.png}{官方 slide 9：按历史准确率方差排序的 one-example RLVR}
\lecturefigure{slides-images/slide-010.png}{官方 slide 10：Exploration 与 extrapolation 的 RL data}
\lecturefigure{slides-images/slide-011.png}{官方 slide 11：High-quality data mix 优先于 parameter alchemy}
\lecturefigure{slides-images/slide-012.png}{官方 slide 12：Data synthesis 的实验结果}
\lecturefigure{slides-images/slide-013.png}{官方 slide 13：可构建 repository 与可验证任务的 Agentic synthesis}

\begin{knowledgebox}{术语消化：RLVR、Rubric 与 Curriculum}
\textbf{RLVR} 是 Reinforcement Learning with Verifiable Rewards，奖励由测试、规则或 ground truth 自动计算；\textbf{rubric} 把开放质量拆成可评分维度；\textbf{curriculum} 根据当前策略能力安排任务难度和顺序。一个样本若在训练过程中成功率变化大，可能比始终全对/全错的样本提供更多区分信号。
\end{knowledgebox}

\begin{warningbox}{合成数据的自洽陷阱}
生成器、环境和 grader 若来自同一模型，可能共同认可错误假设。repository 必须真实 build，任务要由独立测试验证，rubric 要用人工/产品约束校准。数量扩张不能替代 provenance、dedup、难度与 verifier false-positive audit。
\end{warningbox}

\textbf{老师强调：}数据质量和 mix 是日常工作的核心。这里的“只用一个样本”不是建议缩减数据，而是说明 selection 与训练动态可以让少量高信息样本产生可观察更新；真实产品仍需要多环境、多工具和失败模式覆盖。

\subsection{Grader：产品约束如何变成可学习奖励}

slides 15--21 展开 grader 体系。若知道评什么、怎样评，就解决了问题的一半。SWE 产品 grader 可能同时包含 unit test、patch scope、format、behavior、length 和 human-readable summary。多个 grader 的尺度与依赖不同，早期 policy 很难同时通过；可用 relative ranking、curriculum 和 dependency-aware ordering。模型会修改测试或利用漏洞，因此 anti-hacking 是常态工作。

\lecturefigure{slides-images/slide-015.png}{官方 slide 15：Grader/Eval 是 RL experience 的反馈}
\lecturefigure{slides-images/slide-016.png}{官方 slide 16：SWE product grader 的多重约束}
\lecturefigure{slides-images/slide-017.png}{官方 slide 17：Unit test、verifier、LLM grader 与 rubric}
\lecturefigure{slides-images/slide-018.png}{官方 slide 18：Relative ranking、curriculum 与 grader dependency}
\lecturefigure{slides-images/slide-019.png}{官方 slide 19：模型修改测试文件等 reward hacking}
\lecturefigure{slides-images/slide-020.png}{官方 slide 20：独立模型监管、final answer 与 hack 检查}
\lecturefigure{slides-images/slide-021.png}{官方 slide 21：Data/grader 与 efficiency 各占约一半日常工作}

\begin{knowledgebox}{读图：多 grader 不是简单加权平均}
硬约束如 build/test failure 应先作为 gate；软约束如代码风格和解释可在通过后排序。若所有 grader 同时稀疏，早期策略几乎没有正奖励；curriculum 可先训练基础可执行性，再逐步加入产品要求。Relative ranking 比绝对阈值更容易在早期提供差异信号，但必须防止模型优化相对捷径。
\end{knowledgebox}

\begin{warningbox}{Reward hacking 是 grader 的测试用例}
模型修改 test、删除失败文件或输出伪造日志，说明 reward 与真实目标不一致。防御包括保护 hidden tests、只读挂载、diff/行为检查、独立监管模型和人工坏例回灌。\textbf{课堂提示：}讲者说 cheating is normal，是要求团队预设对抗，而不是将其归因于模型道德。
\end{warningbox}

\subsection{Efficiency：Sampling 主导成本，Async RL 解耦等待}

slides 22--24 说明 RL 尤其昂贵，因为模型必须采样探索并等待环境。同步架构中，一个慢 rollout 会阻塞 trainer；异步 RL 让 trainer、多个 environment roller 和 grader 解耦，通过队列持续生产经验。收益来自隐藏环境等待和扩大并发，但也引入 policy staleness、负载不均与版本一致性问题。

\lecturefigure{slides-images/slide-022.png}{官方 slide 22：Sampling 是 Agent RL 的主要效率瓶颈}
\lecturefigure{slides-images/slide-023.png}{官方 slide 23：Trainer 与 rollers 的 asynchronous RL}
\lecturefigure{slides-images/slide-024.png}{官方 slide 24：Environment rollers、graders 与 waiting sync}

\begin{knowledgebox}{术语消化：Roller、Staleness 与 Backpressure}
\textbf{roller} 运行 policy 与环境，生成 trajectory；\textbf{policy staleness} 指 rollout 使用的参数落后于当前 trainer；\textbf{backpressure} 是下游处理不过来时限制上游生产。异步系统要在吞吐与 on-policy 程度之间取舍，并按环境延迟、轨迹长度和 grader 成本做动态调度。
\end{knowledgebox}

\begin{importantbox}{效率指标}
除 GPU utilization，还应记录 valid trajectories/hour、successful learning samples/\$、environment wait ratio、grader latency、discarded stale rollout 和 length distribution。设备忙不代表获得了有效经验；错误/重复轨迹同样会让利用率看起来很高。
\end{importantbox}

\subsection{Exploration、Tool Generalization 与 Solution Length}

slides 26--30 把探索落到 instruction following、工具泛化和长度。策略先要可靠遵循指令，探索才是有意义的；robust IF 使模型适应新工具，MCP 扩大 tool ecosystem。trajectory 越长，sampling 成本越高且错误累积，训练可 cap long rollout，评估则同时看 pass-rate 与长度。百万 token context 也不能解决用户长期不结束会话和无界历史的问题。

\lecturefigure{slides-images/slide-026.png}{官方 slide 26：Instruction following 是有效探索前提}
\lecturefigure{slides-images/slide-027.png}{官方 slide 27：Tool generalization 与 MCP 支持}
\lecturefigure{slides-images/slide-028.png}{官方 slide 28：Solution length 决定 sampling efficiency}
\lecturefigure{slides-images/slide-029.png}{官方 slide 29：训练 cap rollout、评估联合 pass-rate 与长度}
\lecturefigure{slides-images/slide-030.png}{官方 slide 30：Long context 仍无法处理无界会话}

\begin{knowledgebox}{读图：长度是质量与成本的联合变量}
短轨迹可能高效，也可能跳过必要验证；长轨迹可能探索充分，也可能循环。应比较成功条件下的长度、失败长度、工具调用和恢复次数，并对任务难度归一化。训练 cap 是预算控制，不能在截断后把未完成轨迹当普通失败而不区分原因。
\end{knowledgebox}

\begin{warningbox}{Context window 不等于 memory policy}
即使模型可读百万 token，把全部历史塞入上下文仍有成本、噪声、隐私和关注衰减。生产 Agent 需要摘要、结构化状态、检索、过期和用户会话边界。更长窗口扩大上限，不替代记忆治理。
\end{warningbox}

\subsection{LoRA 与统一训练视角}

slides 31--32 提出 LoRA 在 RL 中不仅省显存，还可能通过低秩更新形成 regularization，减少过拟合并改善 generalization。随后将 pre-training、post-training 和 RL 统一：预训练找到通用积木与高效数据，post-training 教模型遵循目标，RL 让模型在环境反馈中探索。三者是能力形成的连续链，而不是互相替代。

\lecturefigure{slides-images/slide-031.png}{官方 slide 31：LoRA 在 RL 中的 regularization 与 generalization}
\lecturefigure{slides-images/slide-032.png}{官方 slide 32：Pre-training、Post-training 与 RL 的分工}

\begin{knowledgebox}{LoRA 的边界}
LoRA 将权重更新限制在低秩子空间，可降低通信/存储并抑制大幅漂移，但 rank 太低会限制能力，adapter merge 与多任务组合也有成本。slide 支持其潜在 regularization 作用，不证明 LoRA 在所有 RL 任务都优于 full fine-tuning；需要同预算、同数据的消融。
\end{knowledgebox}

\subsection{Look Ahead：Testing-time、Training-time 与 Infra Reality}

slides 34--36 展望两类 scaling。testing-time 可让更大/更深模型、并行搜索和 multi-agent interaction 扩大执行；training-time 受限于数据，需 synthesis、self-play 与持续反馈。最后一页回到现实：训练 job 会失败，工程师凌晨查看任务，GPU 资源珍贵，基础设施不稳定是日常问题。

\lecturefigure{slides-images/slide-034.png}{官方 slide 34：Testing-time scaling 的深度、宽度与交互}
\lecturefigure{slides-images/slide-035.png}{官方 slide 35：Training-time scaling、data synthesis 与 self-play}
\lecturefigure{slides-images/slide-036.png}{官方 slide 36：Training model 的作业、GPU 与基础设施现实}

\begin{importantbox}{Scaling 的共同约束}
测试时扩展增加单任务成本，训练时扩展增加数据和系统压力。两者都必须由 grader 证明收益、由长度/预算约束成本、由 infra 支撑并发。若 evaluation 不可靠，更多 compute 只会更快优化错误目标。
\end{importantbox}

\textbf{实践经验：}讲者用“凌晨三点看 job”强调模型训练是持续运营，不是一次脚本。checkpoint、自动恢复、健康检查、环境隔离和可观测性应在算法实验前准备，否则研究结论会被系统故障和不可复现运行污染。

